\section*{Chapitre \ref{ch:b3:complex-kinetics} --- Cinétique complexe~: chaînes, enzymes et oscillations}

\begin{solution}{exo:b3:complex-kinetics:1}
$\ce{Cl2 -> 2 Cl}$~: amorçage~; $\ce{Cl + H2 -> HCl + H}$ et $\ce{H + Cl2 -> HCl + Cl}$~:
propagation~; $\ce{2 Cl -> Cl2}$~: rupture. Porteurs~: Cl et H. Bilan
(somme des deux étapes de propagation)~: $\ce{H2 + Cl2 -> 2 HCl}$.
\end{solution}

\begin{solution}{exo:b3:complex-kinetics:2}
$[\mathrm M]_{1/2} = k_2/k_{-1} = \qty{1.0e-3}{mol/L} = \qty{1.0}{mol/m^3}$~; $p = cRT = 1.0 \times 8.314 \times 750 =
\qty{6.2}{kPa}$.
\end{solution}

\begin{solution}{exo:b3:complex-kinetics:3}
$v/V_{\max} = [\mathrm S]/(K_M + [\mathrm S])$~: $\frac12$, $\frac34$, $\frac{9}{10}$. Pour 90\,\% de $V_{\max}$, il faut $[\mathrm S] = 9K_M$.
\end{solution}

\begin{solution}{exo:b3:complex-kinetics:4}
$k_{\mathrm{cat}}/K_M = 100/\num{0.80e-3} = \qty{1.25e5}{L.mol^{-1}.s^{-1}}$~:
quelque 60\,000 fois au-dessous de la limite de diffusion, et proche de
la valeur typique.
\end{solution}

\begin{solution}{exo:b3:complex-kinetics:5}
AEQS pour \ce{CH3CO}~: $k_2[\ce{CH3}][\ce{CH3CHO}] = k_3[\ce{CH3CO}]$. AEQS pour \ce{CH3}~:
$k_1[\ce{CH3CHO}] - k_2[\ce{CH3}][\ce{CH3CHO}] + k_3[\ce{CH3CO}] - 2k_4[\ce{CH3}]^2 = 0$. En additionnant~: $k_1[\ce{CH3CHO}] =
2k_4[\ce{CH3}]^2$, donc la concentration $[\ce{CH3}] = (k_1/2k_4)^{1/2}[\ce{CH3CHO}]^{1/2}$, et $\dd[\ce{CH4}]/\dd t = k_2[\ce{CH3}][\ce{CH3CHO}] =
k_2(k_1/2k_4)^{1/2}[\ce{CH3CHO}]^{3/2}$.
\end{solution}

\begin{solution}{exo:b3:complex-kinetics:6}
À mi-conversion, $[\ce{H2}] = [\ce{Br2}] = \frac12$, $[\ce{HBr}] = 1$
(unités relatives)~:
$v/v_0 = \frac12 \times (\frac12)^{1/2}/(1 + 0.10 \times 1/0.5) = 0.354/1.2 = 0.29$.
Sans inhibition, ce serait
0.354~: l'inhibition divise encore la vitesse par 1.2.
\end{solution}

\begin{solution}{exo:b3:complex-kinetics:7}
Explosion quand $1000p > 6000/p + 10p^2$, c'est-à-dire $10p^3 - 1000p^2 + 6000 < 0$,
dont les racines positives
valent 2.48 et \qty{99.9}{kPa}~: le mélange explose entre environ 2.5 et
\qty{100}{kPa} (première et deuxième limites du modèle).
\end{solution}

\begin{solution}{exo:b3:complex-kinetics:8}
$1/v = 1/V_{\max} + (K_M/V_{\max})/[\mathrm S]$~: $5.0 = 1/V_{\max} + 2.0K_M/V_{\max}$ et $2.5 = 1/V_{\max} + 0.5K_M/V_{\max}$.
En retranchant~: $K_M/V_{\max} = \qty{1.667}{s.mM.\micro M^{-1}}$, puis $1/V_{\max} = 1.667$~: $V_{\max} = \qty{0.60}{\micro M.s^{-1}}$,
$K_M = \qty{1.0}{mM}$.
\end{solution}

\begin{solution}{exo:b3:complex-kinetics:9}
Même ordonnée à l'origine sur l'axe des $1/v$~: \omterm{def:b3:complex-kinetics:inhibition}{inhibition compétitive}.
$K_M^{\mathrm{app}} = 1/0.417 = \qty{2.40}{mM} = 3K_M$, donc
$1 + 50/K_i = 3$ et $K_i = \qty{25}{\micro M}$.
\end{solution}

\begin{solution}{exo:b3:complex-kinetics:10}
$N = \qty{1.01}{mol/L}$~: $t_{1/2} = \ln(1.01/0.010 - 1)/(0.50 \times 1.01) = \ln100/0.505 = \qty{9.1}{s}$.
La vitesse $kx(N - x)$ est maximale pour $x = N/2$~: c'est le même
instant, le point d'inflexion de la courbe en S.
\end{solution}

\begin{solution}{exo:b3:complex-kinetics:11}
Trace $B - 2$, déterminant 1. $B = 1.5$~: $\lambda = -0.25 \pm 0.968\iu$,
un foyer stable (oscillations amorties vers l'état stationnaire).
$B = 2.5$~: $\lambda = 0.25 \pm 0.968\iu$, un foyer instable~: les
oscillations croissent jusqu'à atteindre le \omterm{def:b3:complex-kinetics:oscillating}{cycle limite}.
\end{solution}

\begin{solution}{exo:b3:complex-kinetics:12}
$1/\tau = \num{1.0e8} \times \num{5.40e-5} + \num{1.0e3} = \qty{6.4e3}{s^{-1}}$, $\tau = \qty{156}{\micro s}$.
Mesurer $\tau$ à plusieurs
concentrations totales~: $1/\tau$ en fonction de $[\mathrm A]_e + [\mathrm B]_e$ est une droite de pente $k_1$ et
d'ordonnée à l'origine $k_{-1}$.
\end{solution}

\begin{solution}{pb:b3:complex-kinetics:1}
\textbf{1.} Au départ, $[\mathrm S]$ est connue et le produit, qui
pourrait inhiber l'\omterm{def:b3:complex-kinetics:enzyme}{enzyme} ou réagir en retour, est absent.
\textbf{2.} Les points de faible $[\mathrm S]$, à grands $1/[\mathrm S]$ et
$1/v$~: ils fixent la pente, et l'inversion amplifie leurs erreurs.
\textbf{3.} Il pondère chaque vitesse telle qu'elle est mesurée~; la
droite de Lineweaver--Burk donne des paramètres biaisés.
\textbf{4.} $K_M$~: $0.773$ est à 1.2 écart-type au-dessous de 0.801~;
$V_{\max}$~: $0.491$ est à 2.2 écarts-types au-dessous de 0.502. Les
deux sont trop bas, comme le laisse attendre le biais.
\textbf{5.} $k_{\mathrm{cat}} = 0.502/0.0050 = \qty{100}{s^{-1}}$.
\textbf{6.} $100/\num{0.801e-3} = \qty{1.25e5}{L.mol^{-1}.s^{-1}}$.
\textbf{7.} Quelque 60\,000 fois au-dessous de la limite de diffusion~;
une \omterm{def:b3:complex-kinetics:enzyme}{enzyme} typique.
\textbf{8.} Compétitive~: $V_{\max}$ inchangée, $K_M$ apparente croissant
avec $[\mathrm I]$.
\textbf{9.} $K_M^{\mathrm{app}} = K_M(1 + [\mathrm I]/K_i) = K_M + (K_M/K_i)[\mathrm I]$.
\textbf{10.} $K_i = 0.799/0.0321 = \qty{24.9}{\micro M}$.
\textbf{11.} $24.9 \pm 4.30 \times 0.9 = 24.9 \pm \qty{3.9}{\micro M}$.
\textbf{12.} La constante de dissociation du complexe enzyme--inhibiteur~:
la concentration d'inhibiteur qui occupe la moitié de l'\omterm{def:b3:complex-kinetics:enzyme}{enzyme} libre.
\textbf{13.} $v_i/v_0 = (K_M + [\mathrm S])/(K_M(1 + [\mathrm I]/K_i) + [\mathrm S])$.
\textbf{14.} $2/3 = 0.67$.
\textbf{15.} $2/12 = 0.17$.
\textbf{16.} $11/21 = 0.52$~: à forte concentration, le substrat l'emporte
sur l'inhibiteur.
\textbf{17.} Pour $[\mathrm S] = K_M$, $v_i/v_0 = 2/(2 + [\mathrm I]/K_i) = 0.10$ donne $[\mathrm I]/K_i = 18$.
\textbf{18.} $18 \times 24.9 = \qty{448}{\micro M}$.
\textbf{19.} La première étape est rapide~: chaque molécule d'\omterm{def:b3:complex-kinetics:enzyme}{enzyme}
libère vite un $\mathrm P_1$
et se trouve piégée sous forme d'\omterm{def:b3:complex-kinetics:enzyme}{acyl-enzyme}~; ensuite, $\mathrm P_1$
n'apparaît qu'au rythme auquel la lente
seconde étape libère l'\omterm{def:b3:complex-kinetics:enzyme}{enzyme}.
\textbf{20.} $1.28/2.0 = 0.64 = (k_2/(k_2 + k_3))^2$, donc $k_2/(k_2 + k_3) = 0.80$.
\textbf{21.} $200/2.0 = \qty{100}{s^{-1}}$, la même valeur que dans la
partie I.
\textbf{22.} $k_2 = 0.80 \times 625 = \qty{500}{s^{-1}}$, $k_3 = \qty{125}{s^{-1}}$.
\textbf{23.} $500 \times 125/625 = \qty{100}{s^{-1}}$.
\textbf{24.} \textbf{$[\mathrm I]_{90\,\%} = 18K_i \approx \qty{0.45}{mM}$
pour $[\mathrm S] = K_M$.}
\end{solution}
